\section*{Summary}
\addcontentsline{toc}{section}{Summary}

\begin{itemize}
  \item \code{Frame} is an \textbf{aggregate}: no constructor, no destructor, no methods. That
    keeps aggregate initialisation readable and makes the class pure data.
  \item \code{MaxDataLen} is \code{static constexpr} --- a type, a scope, visible in the debugger,
    and the same number as the size of the array.
  \item \code{isValid()} is a free function, and the rule is expressed \textbf{once}. The hardware
    validates nothing, so the driver has to --- and only it.
  \item You write no tests in the course. You run them, read them and get them to pass.
  \item Testing decides the course because the hardware neither exists nor is yours until second
    to last.
  \item Read a failing test case in the order test name $\rightarrow$ Arrange $\rightarrow$
    Assert, and open your own code last.
  \item The test suite is more often a more precise requirements specification than the prose.
\end{itemize}
